\documentclass{report}
\usepackage{amsmath,amssymb}
\setlength{\parindent}{0mm}
\setlength{\parskip}{1em}
\begin{document}
\begin{center}
\rule{6in}{1pt} \
{\large Alex Breuer \\
{\bf Heterogeneous resolution of commute time embedding features in Krylov subspaces}}

ATTN: RDRL-CI-N \\ Bldg 120 \\ Aberdeen Proving Ground \\ MD 21005
\\
{\tt alexander.m.breuer.civ@mail.mil}\end{center}

Computing the eigenvectors needed for a commute-time embedding using the
truncated SVD is expensive when the problem is nontrivial. One may use
approximate eigenvectors instead, and significant computational savings
may result; in particular Krylov subspaces have been suggested for this
purpose. One may witness that global features of the graph emerge before
local detail is resolved. We show that this heterogeneous resolution of
graph features is an innate property of Krylov subspace approximations.
Error of global structure shrinks at an asymptotically faster rate than
the error of local detail. We also notice that a similar behavior results
for other Krylov subspace approximations of dimension reduction methods
that use eigenvector projections.


\end{document}
